\phantomsection\label{sec:pub}
\section*{论文发表}
粗体表示本简历作者。星号（*）表示同等贡献。

\subsection*{主要论文（第一作者、共同第一作者或主要学生作者）}
以下论文由我担任第一作者、共同第一作者或主要学生作者。
在每项工作中，我对问题定义、算法设计、论文撰写和修改、代码实现及实验评估中的部分或全部环节作出了重要贡献。

\pubentry{Advanced Materials, 2026}{\href{https://doi.org/10.1002/adma.202522493}{论文} \,/\, \href{https://arxiv.org/abs/2509.13916}{预印本} \,/\, \href{https://github.com/Logan-Lin/AMDEN-code}{代码}}
{Inverse Design of Amorphous Materials with Targeted Properties}
{Jonas A. Finkler*, \me*, Tao Du, Jilin Hu, Morten M. Smedskjaer}
{生成符合目标性能的玻璃等无序材料原子结构，并将其优化为稳定的低能构型。该工作还发布了新的非晶材料数据集，以支持此类设计研究。}

\pubentry{IJCAI, 2026}{\href{https://arxiv.org/abs/2603.29837}{预印本} \,/\, \href{https://github.com/Logan-Lin/DiSGMM-code}{代码}}
{DiSGMM: A Method for Time-varying Microscopic Weight Completion on Road Networks}
{\me, Jilin Hu, Shengnan Guo, Christian S. Jensen, Youfang Lin, Huaiyu Wan}
{在路段之间和单个路段内部的观测均较稀疏时，补全路网中细粒度且随时间变化的交通状态，如特定时段各路段的行驶速度。该方法估计每个路段和时段可能出现的交通状态的完整范围，而非单一数值。}

\pubentry{IJCAI, 2026}{}
{TrajAR: Long-Term Trajectory Prediction at Urban Intersections via Multi-scale Interaction Perception}
{Letian Gong*, \me*, Xinyue Zhang, YiWei Shuang, Guanyu Yao, Cheng Long, Shengnan Guo, Youfang Lin, Huaiyu Wan}
{}

\pubentry{KDD, 2026}{\href{https://arxiv.org/abs/2512.21685}{预印本}}
{RIPCN: A Road Impedance Principal Component Network for Probabilistic Traffic Flow Forecasting}
{Haochen Lv*, \me*, Shengnan Guo, Xiaowei Mao, Hong Nie, Letian Gong, Youfang Lin, Huaiyu Wan}
{预测路网未来交通流量并估计每项预测的不确定性，将交通理论与学习模型结合，捕捉拥堵如何随时间在相连道路之间转移交通流。}

\pubentry{NeurIPS, 2025}{\href{https://neurips.cc/virtual/2025/poster/115959}{海报}}
{TrajMamba: An Efficient and Semantic-rich Vehicle Trajectory Pre-training Model}
{Yichen Liu*, \me*, Shengnan Guo, Zeyu Zhou, Youfang Lin, Huaiyu Wan}
{学习丰富且紧凑的车辆 GPS 轨迹表示，同时捕捉车辆如何移动和为何出行。该方法去除冗余轨迹点，并加入道路和周边地点提供的出行目的信息，从而提升下游任务的精度和效率。}

\pubentry{NeurIPS, 2025}{\href{https://neurips.cc/virtual/2025/oral/117509}{口头报告} \,/\, \href{https://neurips.cc/virtual/2025/poster/117508}{海报} \,/\, \href{https://arxiv.org/abs/2505.12672}{预印本}}
{TransferTraj: A Vehicle Trajectory Learning Model for Region and Task Transferability}
{Tonglong Wei*, \me*, Zeyu Zhou, Haomin Wen, Jilin Hu, Shengnan Guo, Youfang Lin, Gao Cong, Huaiyu Wan}
{从车辆 GPS 轨迹中学习可在不同区域和不同预测任务之间迁移的知识，无须重新训练或维护多个专用模型。该方法将各项任务统一转化为恢复轨迹缺失部分的问题，使单一模型在数据有限时也能处理多种任务。}

\pubentry{NeurIPS, 2025}{\href{https://neurips.cc/virtual/2025/poster/117945}{海报} \,/\, \href{https://arxiv.org/abs/2410.14281}{预印本}}
{PLMTrajRec: A Scalable and Generalizable Trajectory Recovery Method with Pre-trained Language Models}
{Tonglong Wei*, \me*, Youfang Lin, Shengnan Guo, Jilin Hu, Haitao Yuan, Gao Cong, Huaiyu Wan}
{恢复稀疏移动轨迹中的缺失点以重建详细路径，仅需少量稠密训练数据，并能泛化到不同采样率记录的轨迹。}

\pubentry{IEEE TKDE, 2025}{\href{https://ieeexplore.ieee.org/document/11004614}{论文} \,/\, \href{https://arxiv.org/abs/2402.07232}{预印本} \,/\, \href{https://github.com/Logan-Lin/UVTM}{代码}}
{UVTM: Universal Vehicle Trajectory Modeling with ST Feature Domain Generation}
{\me, Jilin Hu, Shengnan Guo, Bin Yang, Christian S. Jensen, Youfang Lin, Huaiyu Wan}
{用单一模型处理出行时间估计、轨迹恢复和轨迹预测等多种车辆 GPS 轨迹任务，无须为每项任务维护独立模型。该方法通过将不完整轨迹重建为稠密且完整的轨迹，在轨迹稀疏或仅有部分特征时仍能保持较高精度。}

\pubentry{IJCAI, 2025}{\href{https://www.ijcai.org/proceedings/2025/411}{论文} \,/\, \href{https://arxiv.org/abs/2405.12459}{预印本} \,/\, \href{https://github.com/Zeru19/PLM4Traj}{代码}}
{TrajCogn: Leveraging LLMs for Cognizing Movement Patterns and Travel Purposes from Trajectories}
{Zeyu Zhou*, \me*, Haomin Wen, Shengnan Guo, Jilin Hu, Youfang Lin, Huaiyu Wan}
{对大语言模型进行适配，使其理解移动轨迹、识别移动方式并推断出行目的。单一模型可以处理不同场景下的多种轨迹分析任务。}

\pubentry{IEEE TKDE, 2025}{\href{https://ieeexplore.ieee.org/document/10818577}{论文} \,/\, \href{https://arxiv.org/abs/2407.12550}{预印本} \,/\, \href{https://github.com/Logan-Lin/UniTE}{代码}}
{UniTE: A Survey and Unified Pipeline for Pre-training Spatiotemporal Trajectory Embeddings}
{\me, Zeyu Zhou, Yicheng Liu, Haochen Lv, Haomin Wen, Tianyi Li, Yushuai Li, Christian S. Jensen, Shengnan Guo, Youfang Lin, Huaiyu Wan}
{综述可复用移动轨迹表示的学习方法，并用共享代码将这些方法整合为统一的模块化框架，使新方法可以在同一基础上构建、比较和评估。}

\pubentry{WWW, 2025}{\href{https://openreview.net/forum?id=KmMSQS6tFn}{论文} \,/\, \href{https://github.com/decisionintelligence/Path-LLM}{代码}}
{Path-LLM: A Multi-Modal Path Representation Learning by Aligning and Fusing with Large Language Models}
{Yongfu Wei*, \me*, Hongfan Gao, Ronghui Xu, Sean Bin Yang, Jilin Hu}
{结合路网结构与描述物理环境和区域背景的文本，学习路网中的路径表示，补充了以往方法忽略的上下文信息。融合后的表示提升了路径排序和出行时间估计等任务的精度，包括少量或无标注数据的场景。}

\pubentry{AAAI, 2025}{\href{https://ojs.aaai.org/index.php/AAAI/article/view/33350}{论文} \,/\, \href{https://arxiv.org/abs/2408.12809}{预印本}}
{DutyTTE: Deciphering Uncertainty in Origin-Destination Travel Time Estimation}
{Xiaowei Mao*, \me*, Shengnan Guo, Yubin Chen, Xingyu Xian, Haomin Wen, Qisen Xu, Youfang Lin, Huaiyu Wan}
{估计起点与终点之间的出行时间并量化每次预测的不确定性。该方法先推断可能采用的路线，再生成可靠的置信区间，以反映不同条件下各路段对整体不确定性的影响。}

\pubentry{NeurIPS, 2024}{\href{https://openreview.net/forum?id=0feJEykDRx}{论文} \,/\, \href{https://neurips.cc/virtual/2024/poster/96914}{海报}}
{Mobility-LLM: Learning Visiting Intentions and Travel Preference from Human Mobility Data with Large Language Models}
{Letian Gong*, \me*, Xinyue Zhang, Yiwen Lu, Xuedi Han, Yichen Liu, Shengnan Guo, Youfang Lin, Huaiyu Wan}
{分析访问地点序列等人员移动数据，学习人们访问地点的原因和出行偏好，并用大语言模型支持未来地点预测等任务。}

\pubentry{SIGMOD, 2024}{\href{https://dl.acm.org/doi/10.1145/3617337}{论文} \,/\, \href{https://arxiv.org/abs/2307.03048}{预印本} \,/\, \href{https://github.com/Logan-Lin/DOT}{代码}}
{Origin-Destination Travel Time Oracle for Map-based Services}
{\me, Huaiyu Wan, Jilin Hu, Shengnan Guo, Bin Yang, Christian S. Jensen, Youfang Lin}
{从连接相同起终点的大量历史行程中学习，估计给定出发时间的起终点出行时间。该方法先推断两点之间的可能路线，再预测沿线出行时间，从而提升地图导航服务的估计精度。}

\pubentry{IEEE TKDE, 2023}{\href{https://ieeexplore.ieee.org/abstract/document/10375102}{论文} \,/\, \href{https://arxiv.org/abs/2207.14539}{预印本} \,/\, \href{https://github.com/Logan-Lin/MMTEC}{代码}}
{Pre-training General Trajectory Embeddings with Maximum Multi-view Entropy Coding}
{\me, Huaiyu Wan, Shengnan Guo, Jilin Hu, Christian S. Jensen, Youfang Lin}
{从无标注数据中学习通用移动轨迹表示，同时捕捉出行行为以及空间和时间模式。所得表示避免针对特定任务的偏差，因而能够迁移到多种下游任务。}

\pubentry{IEEE TKDE, 2022}{\href{https://ieeexplore.ieee.org/abstract/document/9351627}{论文} \,/\, \href{https://github.com/Logan-Lin/TALE}{代码}}
{Pre-training Time-Aware Location Embeddings from Spatial-Temporal Trajectories}
{Huaiyu Wan, \me, Shengnan Guo, Youfang Lin}
{从时空轨迹中学习地点向量表示，既捕捉地点的位置，也捕捉人们何时访问地点，因为访问时间能够反映地点用途。这些时间感知表示不依赖标注数据完成预训练，并提升多种基于地点的下游预测任务的性能。}

\pubentry{AAAI, 2021}{\href{https://ojs.aaai.org/index.php/AAAI/article/view/16548}{论文} \,/\, \href{https://github.com/Logan-Lin/CTLE}{代码}}
{Pre-training Context and Time Aware Location Embeddings from Spatial-Temporal Trajectories for User Next Location Prediction}
{\me, Huaiyu Wan, Shengnan Guo, Youfang Lin}
{从移动轨迹中预训练地点表示，捕捉地点含义随周边环境和访问时间产生的变化，从而更准确地预测用户的下一个地点。}

\subsection*{其他论文（参与作者）}
以下论文由我担任参与作者。
在这些工作中，我主要通过技术讨论、算法改进、实验或写作作出贡献。

\pubentry{KDD, 2026}{}
{Low-Rank Prior-Induced Consistency Flow Matching for Efficient Traffic Imputation}
{Xiaowei Mao, Tingrui Wu, Yang Yawen, Shengnan Guo, \me, Shilong Zhao, Haochen Lv, Youfang Lin, Huaiyu Wan}
{}

\pubentry{IJCAI, 2026}{}
{G-VTM: A Multimodal Vision-Trajectory Model for Generalized Vehicle Trajectory Prediction}
{Xinyue Zhang, Letian Gong, \me, Jinjun Cheng, Junlin Zhang, Guanyu Yao, Shengnan Guo, Youfang Lin, Shaojiang Wang, Huaiyu Wan}
{}

\pubentry{DASFAA, 2026}{\href{https://link.springer.com/chapter/10.1007/978-981-92-0375-8_24}{论文}}
{ForceTraj: Modeling Realistic Intention and Heterogeneous Interactions for Multi-modal Trajectory Prediction}
{Guanyu Yao, \me, Letian Gong, Xinyue Zhang, Yiwei Shuang, Shengnan Guo, Huaiyu Wan}
{通过推断变道等驾驶意图，并建模周边车辆之间的相互影响，预测自动驾驶场景中的车辆未来路径。该方法生成多条可能的未来轨迹，而非单一结果，从而提升真实驾驶数据集上的预测精度。}

\pubentry{KDD, 2026}{\href{https://arxiv.org/abs/2509.00053}{预印本}}
{Traj-MLLM: Can Multimodal Large Language Models Reform Trajectory Data Mining?}
{Shuo Liu, Di Yao, \me, Gao Cong, Jingping Bi}
{将原始轨迹转换为保留时空结构的图像与文本组合输入，使多模态大语言模型无须训练即可分析不同区域和任务中的移动轨迹。单一配置可以处理多种轨迹分析任务。}

\pubentry{AAAI, 2026}{\href{https://ojs.aaai.org/index.php/AAAI/article/view/37139}{论文} \,/\, \href{https://arxiv.org/abs/2511.12489}{预印本}}
{SculptDrug: A Spatial Condition-Aware Bayesian Flow Model for Structure-based Drug Design}
{Qingsong Zhong, Haomin Yu, \me, Wangmeng Shen, Long Zeng, Jilin Hu}
{生成适配目标蛋白质三维结构的药物分子，使分子契合目标蛋白质表面，并同时适配其整体形态和局部细节，从而为基于结构的药物发现生成更准确的候选分子。}

\pubentry{AAAI, 2026}{\href{https://ojs.aaai.org/index.php/AAAI/article/view/38581}{论文} \,/\, \href{https://arxiv.org/abs/2601.04572}{预印本}}
{Spatial-Temporal Feedback Diffusion Guidance for Controlled Traffic Imputation}
{Xiaowei Mao, Huihu Ding, \me, Tingrui Wu, Shengnan Guo, Dazhuo Qiu, Feiling Fang, Jilin Hu, Huaiyu Wan}
{通过引导生成过程贴近已有观测值来补全交通数据中的缺失值，并针对不同位置和时间调整生成过程受观测约束的强度，使观测稀疏位置的重建更加准确。}

\pubentry{KDD, 2025}{\href{https://dl.acm.org/doi/10.1145/3690624.3709325}{论文} \,/\, \href{https://arxiv.org/abs/2412.10859}{预印本} \,/\, \href{https://github.com/decisionintelligence/DUET}{代码}}
{DUET: Dual Clustering Enhanced Multivariate Time Series Forecasting}
{Xiangfei Qiu, Xingjian Wu, \me, Chenjuan Guo, Jilin Hu, Bin Yang}
{通过适应时序模式的动态变化，并在降低噪声变量影响的同时建模不同变量之间的关系，提升多变量时间序列预测性能。}

\pubentry{IEEE TKDE, 2024}{\href{https://www.computer.org/csdl/journal/tk/5555/01/10679607/20b3hlbjBOo}{论文} \,/\, \href{https://arxiv.org/abs/2402.07369}{预印本} \,/\, \href{https://github.com/wtl52656/Diff-RNTraj}{代码}}
{Diff-RNTraj: A Structure-aware Diffusion Model for Road Network-constrained Trajectory Generation}
{Tonglong Wei, Youfang Lin, Shengnan Guo, \me, Yiheng Huang, Chenyang Xiang, Yuqing Bai, Menglu Ya, Huaiyu Wan}
{生成遵循真实路网并带有道路层面信息的合成车辆轨迹，缓解隐私问题导致的公开轨迹数据不足。每个生成点均关联到具体路段和移动比例，使轨迹符合实际并可直接用于下游轨迹挖掘任务。}

\pubentry{IEEE TKDE, 2024}{\href{https://ieeexplore.ieee.org/document/10836764}{论文}}
{STCDM: Spatio-Temporal Contrastive Diffusion Model for Check-In Sequence Generation}
{Letian Gong, Shengnan Guo, \me, Yichen Liu, Erwen Zheng, Yiwei Shuang, Youfang Lin, Jilin Hu, Huaiyu Wan}
{生成逼真的合成签到序列，捕捉人员移动及地点访问的时空模式，使研究人员无须依赖因隐私问题而稀缺的真实记录，也能研究和使用此类数据。}

\pubentry{IEEE TKDE, 2024}{\href{https://www.computer.org/csdl/journal/tk/5555/01/10517676/1WCj0j0FljW}{论文} \,/\, \href{https://arxiv.org/abs/2404.19141}{预印本} \,/\, \href{https://github.com/wtl52656/MM-STGED}{代码}}
{Micro-Macro Spatial-Temporal Graph-based Encoder-Decoder for Map-Constrained Trajectory Recovery}
{Tonglong Wei, Youfang Lin, \me, Shengnan Guo, Lan Zhang, Huaiyu Wan}
{重建稀疏 GPS 轨迹中的缺失点，使恢复后的路径保持在道路网络上。该方法结合单次出行的细节与大量行程共享的宏观出行模式，更准确地填补轨迹缺失。}

\pubentry{KBS, 2024}{\href{https://www.sciencedirect.com/science/article/pii/S0950705123010730}{论文} \,/\, \href{https://github.com/wtl52656/IAGCN}{代码}}
{Inductive and Adaptive Graph Convolution Networks Equipped with Constraint Task for Spatial-Temporal Traffic Data Kriging}
{Tonglong Wei, Youfang Lin, Shengnan Guo, \me, Yiji Zhao, Xiyuan Jin, Zhihao Wu, Huaiyu Wan}
{通过从已有传感器中学习，估计未布设传感器位置的交通状态，并能处理新加入的位置，无须重新训练模型。}

\pubentry{IEEE TKDE, 2024}{\href{https://ieeexplore.ieee.org/document/10613424}{论文} \,/\, \href{https://arxiv.org/abs/2407.15899}{预印本}}
{Spatial-Temporal Cross-View Contrastive Pre-Training for Check-in Sequence Representation Learning}
{Letian Gong, Huaiyu Wan, Shengnan Guo, Li Xiucheng, \me, Erwen Zheng, Tianyi Wang, Zeyu Zhou, Youfang Lin}
{从基于位置的服务数据中学习通用用户签到序列表示，同时捕捉用户访问的地点和时间。所得表示能够迁移到多种需要理解用户移动行为的下游任务。}

\pubentry{AAAI, 2023}{\href{https://ojs.aaai.org/index.php/AAAI/article/view/25546}{论文} \,/\, \href{https://github.com/LetianGong/CACSR}{代码}}
{Contrastive Pre-training with Adversarial Perturbations for Check-In Sequence Representation Learning}
{Letian Gong, Youfang Lin, Shengnan Guo, \me, Tianyi Wang, Erwen Zheng, Zeyu Zhou, Huaiyu Wan}
{学习能够捕捉人员访问地点的时空模式的通用签到序列表示，使同一组特征可以迁移到多种移动任务。该方法无须手动设计训练数据变换，所得结果优于特定任务模型且更加稳定。}

\pubentry{ESWA, 2023}{\href{https://www.sciencedirect.com/science/article/pii/S0957417423012241}{论文}}
{Adversarial Self-Attentive Time-Variant Neural Networks for Multi-Step Time Series Forecasting}
{Changxia Gao, Ning Zhang, Youru Li, \me, Huaiyu Wan}
{一次预测时间序列未来多个时间步，并适应序列模式随时间产生的变化。训练过程判断预测结果是否真实并能否从近期输入自然延续，从而保持长期预测精度。}

\pubentry{APIN, 2023}{\href{https://link.springer.com/article/10.1007/s10489-023-05057-7}{论文}}
{Multi-scale Adaptive Attention-based Time-Variant Neural Networks for Multi-step Time Series Forecasting}
{Changxia Gao, Ning Zhang, Youru Li, \me, Huaiyu Wan}
{通过适应时序模式变化并捕捉不同时间尺度的规律，预测时间序列未来多个时间步。该方法在气候和能源消耗数据上的预测精度优于现有方法。}

\pubentry{NeurIPS, 2023}{\href{https://openreview.net/forum?id=y08bkEtNBK}{论文} \,/\, \href{https://github.com/Water2sea/WITRAN}{代码}}
{WITRAN: Water-wave Information Transmission and Recurrent Acceleration Network for Long-range Time Series Forecasting}
{Yuxin Jia, Youfang Lin, Xinyan Hao, \me, Shengnan Guo, Huaiyu Wan}
{通过共同捕捉不同时间尺度的规律及其重复结构，对时间序列进行长期预测。该方法保持较高计算效率，在很长的预测范围内仍兼具速度和精度。}
